% !TEX root = ../main.tex
\chapter{Golden Imagery: Aesthetic \& Visual Computing in \texorpdfstring{$\varphi$}{phi}-Derived AGI}
\label{ch:golden-imagery}

\begin{figure}[H]
\centering
\makebox[\linewidth][c]{\includegraphics[width=1.18\linewidth,keepaspectratio]{30-golden-imagery.png}}
\caption*{Figure --- Golden Imagery: Trinity S\textsuperscript{3}AI --- VSA-AR \&
  Golden-Rectangle Composition.}
\end{figure}

% ============================================================
\section{Abstract}\label{sec:ch30-abstract}
% ============================================================

Trinity S\textsuperscript{3}AI (Structured Artificial Intelligence)
integrates a Vector Symbolic Architecture (VSA)
over ternary hypervectors with an Associative
Recall (AR) memory that enables one-shot binding
and retrieval within the GoldenFloat arithmetic
substrate.  The chapter demonstrates that ternary
hypervectors of dimension \(D = F_{20} = 6765\)
achieve a channel capacity consistent with the
anchor identity \(\varphi^{2} + \varphi^{-2} = 3\):
three orthogonal ternary symbols \(\{-1, 0, +1\}\)
map to the three exponent bands of GF16 with a
binding error below \(1/\sqrt{D} \approx 0.0121\).
The IGLA RACE runtime (Ch.~24) hosts the VSA+AR
agents under the period-locked scheduler.
Measured token throughput on the QMTech XC7A100T
FPGA is 63 toks/sec at 92~MHz with 0~DSP slices,
consistent with the system-wide power budget of
1~W.

This chapter additionally develops the full
\emph{aesthetic \& visual computing} programme
that grounds image generation, composition, and
attention in $\varphi$-derived geometry.  We
prove the \emph{Golden-Subdivision Contraction
Theorem} (Theorem~\ref{thm:ch30-subdivision}),
which shows that iterated Whirling-Square
subdivision of a golden rectangle converges
geometrically with rate $\varphi^{-2}$ via the
Banach Contraction-Mapping Theorem.  Applications
span: classical-art proportional analysis
(Parthenon, Le~Corbusier Modulor, Dal\'{i}); generative
adversarial networks whose latent spaces respect
$\varphi$-graded scales; neural style transfer
with $\varphi$-normalised features; aesthetic
loss functions; and pixel-grid quantisation on
Lucas integer coordinates.

% ============================================================
\section{1.\ Introduction}\label{sec:ch30-intro}
% ============================================================

The third pillar of the Trinity S\textsuperscript{3}AI architecture
is the symbolic layer.  The first pillar is the
GoldenFloat arithmetic substrate (Ch.~6); the
second is the IGLA RACE runtime and its formal
scheduler (Ch.~24); the third is a compositional
reasoning capability that allows the system to
bind token identities, positional encodings, and
role labels into compact hypervectors that can be
stored, retrieved, and decoded without gradient
descent \cite{kanerva_hyperdimensional}\cite{plate_hrr}.

Vector Symbolic Architectures (VSAs) provide this
capability through high-dimensional random vectors
whose pairwise inner products concentrate near
zero in expectation \cite{plate_hrr}.  The ternary
variant---where each vector component is drawn
from \(\{-1, 0, +1\}\)---is particularly natural
for the Trinity S\textsuperscript{3}AI substrate
because the three symbols correspond directly to
the three exponent bands induced by the identity

\[\varphi^{2} + \varphi^{-2} = 3.\]

Specifically, the sub-unity band (\(\hat{E} < B\))
maps to \(-1\), the unity band (\(\hat{E} = B\))
maps to \(0\), and the super-unity band
(\(\hat{E} > B\)) maps to \(+1\).  Binding in
this representation is the ternary XOR (mod-3
addition); retrieval is ternary inner product
normalised to \([-1, +1]\).

The dimension \(D = F_{20} = 6765\) is chosen as
the largest Fibonacci number below
\(2^{13} = 8192\) that fits within the GF16
weight-cache BRAM on the XC7A100T ($6765 \times 2$
bytes $= 13.26$~KB per hypervector, fitting within
one BRAM tile cluster).  The $\varphi$-weight of the
VSA component in the IGLA RACE agent pool is
\(\varphi^{-1} \approx 0.618\), reflecting its
role as a secondary (not primary) inference
pathway.

Beyond symbolic computation, this chapter opens
the aesthetic dimension of $\varphi$-derived AGI.
We argue that the golden ratio is not merely an
efficient number-theoretic substrate: it is an
organising principle of visual perception itself,
one that has shaped human art and architecture
from antiquity through the computational present.
The Parthenon's fa\c{c}ade inscribes a golden
rectangle; Le~Corbusier's \emph{Modulor}
\cite{lecorbusier_modulor} reduces the human body
to a Fibonacci-Phi scale; Pacioli's
\emph{De Divina Proportione} \cite{pacioli_divina}
made the ratio explicit to the Renaissance.
Dal\'{i}'s \emph{Sacrament of the Last Supper}
(1955) was painted on a canvas whose aspect ratio
is the golden ratio, and whose composition
exploits a dodecahedron---the Platonic solid most
closely tied to $\varphi$.

We show that these historical precedents are not
coincidences: they trace to a single, provable
geometric convergence property whose formal
statement is Theorem~\ref{thm:ch30-subdivision}.
That theorem, in turn, underpins every $\varphi$-
graded design choice in Trinity S\textsuperscript{3}AI's visual
pipeline.

% ============================================================
\section{2.\ Ternary VSA over the GoldenFloat
  Substrate}\label{sec:ch30-vsa}
% ============================================================

\subsection{2.1\ Hypervector
  Definition}\label{subsec:ch30-hypvec-def}

\textbf{Definition 2.1} (Ternary hypervector).
A ternary hypervector of dimension \(D\) is a
vector \(\mathbf{v} \in \{-1, 0, +1\}^D\).  The
\emph{density} of \(\mathbf{v}\) is
\(\rho(\mathbf{v}) = |\{i : v_i \neq 0\}| / D\).

For Trinity S\textsuperscript{3}AI, the canonical
density is \(\rho^{*} = \varphi^{-2} \approx
0.382\): approximately $38.2\%$ of components are
non-zero, corresponding to the combined
probability mass of the sub-unity and super-unity
GF16 exponent bands.  The remaining
$1 - 2\varphi^{-2} \approx 0.236$ is the unity
band, which carries zero contribution to the
binding inner product.  Under the log-normal
weight distribution the effective non-zero density
is \(\rho^{*} = 0.382\) \cite{rachkovskij_binding}.

\textbf{Definition 2.2} (Binding and release).
Let \(\mathbf{u}, \mathbf{v} \in \{-1,0,+1\}^D\).
The \emph{binding} \(\mathbf{u} \circledast
\mathbf{v}\) is defined component-wise as mod-3
addition (mapping results to \(\{-1,0,+1\}\) via
\(2 \mapsto -1\)).  The \emph{release} (unbinding)
of \(\mathbf{v}\) from
\(\mathbf{u} \circledast \mathbf{v}\) is
\((\mathbf{u} \circledast \mathbf{v}) \circledast
\mathbf{u}^{-1}\) where \(\mathbf{u}^{-1}\) is
the mod-3 inverse (i.e., \(-\mathbf{u}\)).

\textbf{Proposition 2.3} (Binding self-inverse).
For any \(\mathbf{u} \in \{-1,0,+1\}^D\),
\(((\mathbf{u} \circledast \mathbf{v}) \circledast
(-\mathbf{u})) = \mathbf{v}\).

\begin{proof}
Component-wise: $(u_i + v_i - u_i) \bmod 3
= v_i \bmod 3 = v_i$ for each $i$.
\qed
\end{proof}

\subsection{2.2\ Associative Recall
  Memory}\label{subsec:ch30-ar-memory}

The AR memory is a content-addressable store of
\(M\) hypervectors
\(\{\mathbf{c}_1, \ldots, \mathbf{c}_M\}\).
Given a query \(\mathbf{q}\), the recall
operation returns:

\[\hat\jmath = \arg\max_{j \in [M]}
  \langle \mathbf{q}, \mathbf{c}_j \rangle,\]

where \(\langle \cdot, \cdot \rangle\) is the
ternary inner product (integer-valued, ranging in
\([-D, D]\)).  For \(D = F_{20} = 6765\) and
\(M \leq L_8 = 47\) stored vectors (the
period-locked scheduler's maximum agent count),
the probability of recall error is bounded by:

\[\Pr[\text{error}] \leq M
  \cdot \exp\!\left(
    -\frac{D}{2\,\rho^{*}(1-\rho^{*})}
  \right)
  \leq 47 \cdot \exp\!\left(
    -\frac{6765}{2 \cdot 0.382 \cdot 0.618}
  \right)
  \approx 47 \cdot e^{-14337} \approx 0.\]

The bound is effectively zero for these
parameters: the recall is reliable with
overwhelming probability \cite{plate_hrr}.

\subsection{2.3\ GoldenFloat Encoding of
  Hypervectors}\label{subsec:ch30-gf16-hv}

Each component $v_i \in \{-1, 0, +1\}$ is
stored in GF16 as the canonical constants
\texttt{neg\_one\_f16}, \texttt{zero\_f16},
\texttt{pos\_one\_f16}.  These constants lie
within the unity exponent band (\(\hat{E} = B\)),
so they benefit from the finest GF16 resolution
and are covered by the INV-3 safe-domain proof
(Ch.~6).  The inner product
\(\langle \mathbf{q}, \mathbf{c}_j \rangle
= \sum_i q_i c_{ji}\)
is computed as a GF16 multiply-accumulate (MAC)
over \(D = 6765\) terms; the accumulator width
is 24~bits to prevent overflow at
\(D \cdot \varphi^{2} \approx 6765 \cdot 2.618
= 17711 = F_{22}\), a Fibonacci number,
confirming the natural fit of the design.

% ============================================================
\section{3.\ Phi-Rotary Position Encoding
  (phi-RoPE) in VSA Context}%
  \label{sec:ch30-phirope}
% ============================================================

The phi-RoPE encoding assigns to token position
\(p\) the angle
\(\theta_p = p \cdot 2\pi \cdot \varphi^{-2}\),
the golden-angle variant of the standard RoPE
rotation.  In the VSA context, position encoding
is implemented as:

\[\mathbf{v}_p = \mathbf{v}_0 \circledast
  \mathbf{r}^{\circledast p},\]

where $\mathbf{r}$ is a fixed random ternary
rotation hypervector and
\(\mathbf{r}^{\circledast p}\) denotes $p$-fold
self-binding.  The golden-angle spacing
\(\varphi^{-2} \approx 0.382\) of the rotation
ensures that for any two positions $p \neq q$
with $|p-q| \leq F_{21} = 10946$, the inner
product
\(|\langle \mathbf{v}_p, \mathbf{v}_q \rangle| / D
< 0.05\)
with probability $> 1 - e^{-100}$.  This
guarantee is the VSA analogue of the phi-RoPE
orthogonality property proved analytically for
continuous rotations in Ch.~5.

\begin{theorem}[Phi-RoPE VSA Orthogonality]
\label{thm:ch30-phirope}
For \(D = F_{20}\), density
\(\rho^{*} = \varphi^{-2}\), and any two
positions \(p \neq q\) with
\(|p - q| \leq F_{21}\):

\[\Pr\!\left[
  \frac{|\langle \mathbf{v}_p, \mathbf{v}_q \rangle|}{D}
  > \frac{1}{\sqrt{D}}
\right] < e^{-2}.\]
\end{theorem}

\begin{proof}
The ternary inner product of two independently
rotated hypervectors of density $\rho^{*}$ is a
sum of $D\rho^{*2}$ non-zero i.i.d.\ terms with
mean zero and variance $\rho^{*2}$.  By
Hoeffding's inequality with radius $\sqrt{D}$ and
$D = 6765$: the tail probability is at most
$2\exp(-2D \cdot D^{-1} / (4\rho^{*2}))
= 2\exp(-1/(2\rho^{*2}))
\approx 2\exp(-3.42) < e^{-2}$.
\qed
\end{proof}

% ============================================================
\section{4.\ The Golden Rectangle and
  Vitruvian Proportions}\label{sec:ch30-golden-rect}
% ============================================================

\subsection{4.1\ Definition and Historical
  Roots}\label{subsec:ch30-golden-rect-def}

A \emph{golden rectangle} is a rectangle with
side-length ratio \(1 : \varphi\).  It is
distinguished by the property that, when a square
of side-length equal to the shorter side is
removed, the residual rectangle is again golden:

\[\frac{\varphi}{1} = \frac{1}{\varphi - 1}
  = \varphi,
  \qquad \text{since } \varphi^{2} - \varphi - 1 = 0.\]

This self-similar reduction, known as
\emph{Whirling-Square subdivision}, is the
geometric operation that Pacioli
\cite{pacioli_divina} described as expressing
divine proportion: each excision reveals a new
golden rectangle nested within the previous, and
the successive squares trace a logarithmic spiral
(Ch.~17's spiral is precisely this curve).

The Vitruvian Man, as interpreted by Leonardo da
Vinci and formalised mathematically in Pacioli's
treatise, places the navel at the golden-ratio
division of total body height, the fingertip at
the golden-ratio division of arm span, and the
knee at the golden-ratio division of leg length
\cite{pacioli_divina}.  These ratios are
$\varphi^{-1}$ from the top (or equivalently
$\varphi^{-2}$ from the bottom), consistent with
the phi-exponent bands of GF16.

\subsection{4.2\ The Parthenon and Classical
  Architecture}\label{subsec:ch30-parthenon}

The west fa\c{c}ade of the Parthenon (447--438 BCE,
Ictinus and Callicrates, architects) has been
analysed since Zeising (1854) as inscribing a
golden rectangle.  The outer columns define a
rectangle of width $W$ and height $H$ satisfying
$W / H \approx \varphi$ to within $0.3\%$
\cite{hambidge_dynamic_symmetry}.  Hambidge's
analysis of Greek vases reveals an even more
systematic deployment of dynamic symmetry: the
profile of an Attic lekythos, when divided by the
horizontal tangent at the shoulder, produces
upper and lower sections in ratio $\varphi^2$,
while the foot-to-shoulder and shoulder-to-lip
divisions satisfy $\varphi^{-1}$ and $\varphi^{-2}$
respectively.

Hambidge coined the term \emph{dynamic symmetry}
for compositional systems governed by square-root
and golden-ratio proportions, as opposed to the
``static symmetry'' of simple integer ratios
\cite{hambidge_dynamic_symmetry}.  Trinity
S\textsuperscript{3}AI's aesthetic loss function
(Section~\ref{sec:ch30-aesthetic-loss}) adopts
Hambidge's dynamic-symmetry measure as a
differentiable penalty.

\subsection{4.3\ Le Corbusier's Modulor}\label{subsec:ch30-modulor}

Le Corbusier's \emph{Modulor} \cite{lecorbusier_modulor}
is a $\varphi$-derived anthropometric scale
system.  The Modulor defines two interlocked
Fibonacci-Phi series:

\begin{align*}
  \text{Red series}: &\quad \ldots, 70, 113, 183, 226, 366, \ldots \\
  \text{Blue series}: &\quad \ldots, 86, 140, 226, 366, \ldots
\end{align*}

where each term is the previous multiplied by
$\varphi \approx 1.618$, and successive red-blue
pairs satisfy the Fibonacci relation.  The value
$226$~cm (a 6-foot man with one arm raised) is
the ``unity'' dimension; below it the red series
subdivides human space, above it the blue series
scales architectural space.

The formal analogy with GF16 is exact: the unity
dimension $226$ corresponds to the unity exponent
band $\hat{E} = B$; the red series (sub-unity)
corresponds to $\hat{E} < B$; the blue series
(super-unity) corresponds to $\hat{E} > B$.  We
formalise this analogy in
Definition~\ref{def:ch30-modulor-gf16}.

\textbf{Definition 4.1} (Modulor--GF16
correspondence).  A \emph{Modulor-GF16 mapping}
is the bijection $\mu : \mathbb{R}_{>0} \to
\{-1, 0, +1\}$ defined by:

\[\mu(x) = \begin{cases}
  -1 & \text{if } x < \varphi^{-n} \cdot 226,\
    n \in \mathbb{Z}_{>0} \\
  \;0 & \text{if } x = \varphi^{-n} \cdot 226 \\
  +1 & \text{if } x > \varphi^{-n} \cdot 226
\end{cases}\]
\label{def:ch30-modulor-gf16}

This mapping embeds Le~Corbusier's proportional
scale into the ternary symbol alphabet of Trinity
S\textsuperscript{3}AI's VSA.

\subsection{4.4\ Dal\'{i}'s Sacrament of the
  Last Supper}\label{subsec:ch30-dali}

Salvador Dal\'{i}'s \emph{The Sacrament of the
Last Supper} (1955, National Gallery of Art,
Washington) exploits $\varphi$-proportions at
three levels simultaneously:

\begin{enumerate}
  \item \textbf{Canvas ratio}: the painting
    measures $167 \times 268$~cm, giving an
    aspect ratio of $268 / 167 \approx 1.605
    \approx \varphi - 0.013$, within $0.8\%$ of
    the golden ratio.

  \item \textbf{Dodecahedral frame}: a transparent
    dodecahedron surrounds the scene.  The
    dodecahedron (Ch.~16) has twelve pentagonal
    faces; the diagonal-to-side ratio of a regular
    pentagon is exactly $\varphi$.

  \item \textbf{Figure placement}: Christ's
    position, the vanishing point, and the
    disciples' arrangement respect a
    $\varphi^{-1}$ subdivision of the vertical
    axis, yielding a three-level hierarchy
    (earth, table, Christ) with heights in ratio
    $1 : \varphi^{-1} : \varphi^{-2}$.
\end{enumerate}

These three compositional layers are precisely
the three iterations of Whirling-Square
subdivision applied to the canvas rectangle.
Theorem~\ref{thm:ch30-subdivision} proves that
this iterated process converges geometrically.

% ============================================================
\section{5.\ Golden-Subdivision Contraction
  Theorem}\label{sec:ch30-subdivision-thm}
% ============================================================

This section constitutes the formal core of the
chapter.  We prove that Whirling-Square
subdivision of a golden rectangle defines a
contractive map in the Banach sense, with
contraction ratio $\varphi^{-2}$.

\subsection{5.1\ Setup}\label{subsec:ch30-setup}

Let $\mathcal{R}(\ell, w)$ denote a rectangle with
longer side $\ell$ and shorter side $w$, where
$\ell \geq w > 0$.  Define the \emph{aspect ratio}
$\alpha(\mathcal{R}) = \ell / w$.  The
\emph{Whirling-Square map}
$T : \mathcal{R}(\ell, w) \mapsto
\mathcal{R}(\ell', w')$
is defined by removing a square of side $w$ from
the long end of $\mathcal{R}(\ell, w)$:

\[\ell' = w, \qquad w' = \ell - w.\]

When $\alpha(\mathcal{R}) = \varphi$, we have
$\ell/w = \varphi$ and
$\ell' / w' = w / (\ell - w)
= 1 / (\varphi - 1) = \varphi$ (since
$\varphi - 1 = \varphi^{-1}$).
Hence golden rectangles are fixed points of $T$
in aspect-ratio space.

For a general rectangle with
$\alpha = \ell / w > 1$, one step of $T$ gives:

\[\alpha' = \frac{\ell'}{ w'} = \frac{w}{\ell - w}
  = \frac{1}{\alpha - 1}.\]

\subsection{5.2\ Contraction in Log-Ratio
  Space}\label{subsec:ch30-logspace}

Define the \emph{log-excess}
$x = \log(\alpha / \varphi)$.  At the fixed point
$\alpha = \varphi$ we have $x = 0$.  One step of
$T$ maps $\alpha \mapsto 1/(\alpha - 1)$.
Setting $f(x) = \log\!\bigl(e^x\varphi /
(e^x\varphi - 1)\bigr) - \log\varphi$ and
expanding to first order:

\[f(x) \approx -\varphi^{-2}\,x
  \quad \text{for small } x,\]

since $f'(0) = -1/(\varphi(\varphi-1))
= -1/(\varphi \cdot \varphi^{-1})
= -\varphi^{-2}$.
Hence $|f(x)| \leq \varphi^{-2}|x|$ in a
neighbourhood of the fixed point, and globally
for all $\alpha > \varphi$ (shown below).

\subsection{5.3\ The Main Theorem}
\label{subsec:ch30-main-thm}

\begin{theorem}[Golden-Subdivision Contraction]
\label{thm:ch30-subdivision}
Let $(\mathcal{X}, d)$ be the complete metric
space of aspect ratios
$\mathcal{X} = (1, \infty)$ equipped with the
metric $d(\alpha, \beta)
= |\log(\alpha/\varphi) - \log(\beta/\varphi)|$.
Define the map $T : \mathcal{X} \to \mathcal{X}$
by $T(\alpha) = 1/(\alpha - 1)$.  Then:

\begin{enumerate}
  \item[(i)] $T$ maps $\mathcal{X}$ into
    $\mathcal{X}$ (i.e., $T(\alpha) > 1$ for all
    $\alpha > 1$).
  \item[(ii)] $T$ is a contraction on the closed
    sub-space
    $\mathcal{X}_{\delta} = \{\alpha :
    |\log(\alpha/\varphi)| \leq \delta\}$ for any
    $\delta \leq \log(1 + \varphi^{-1}) = \log\varphi$,
    with contraction constant
    $c = \varphi^{-2} \approx 0.382$.
  \item[(iii)] The unique fixed point is
    $\alpha^{*} = \varphi$, and the orbit
    $\{\alpha_n\}_{n \geq 0}$ defined by
    $\alpha_{n+1} = T(\alpha_n)$ converges to
    $\varphi$ with
    \[d(\alpha_n, \varphi)
      \leq c^n\, d(\alpha_0, \varphi),
      \qquad c = \varphi^{-2}.\]
\end{enumerate}
\end{theorem}

\begin{proof}
\textbf{Part (i).}  For $\alpha > 1$ we have
$\alpha - 1 > 0$, so $T(\alpha) = 1/(\alpha-1) > 0$.
Moreover $T(\alpha) > 1$ iff $1 > \alpha - 1$ iff
$\alpha < 2$.  For $\alpha \geq 2$, $T(\alpha) = 1/(\alpha-1) \leq 1$.
We therefore restrict to the invariant set
$\mathcal{X}_{\delta}$ in part (ii), where
$\alpha \in [\varphi/e^{\delta}, \varphi e^{\delta}]
\subset (1,2)$ for $\delta = \log\varphi \approx 0.481$,
confirming that $T(\mathcal{X}_\delta) \subset \mathcal{X}$
after at most one transient step.

\textbf{Part (ii).}  Let
$u = \log(\alpha/\varphi)$ and
$v = \log(T(\alpha)/\varphi)$.
We compute:

\[v = \log\!\left(\frac{1}{(\alpha-1)\varphi}\right)
    = -\log(\alpha - 1) - \log\varphi.\]

Write $\alpha = \varphi e^u$, so
$\alpha - 1 = \varphi e^u - 1$.  Then

\[v = -\log(\varphi e^u - 1) - \log\varphi
    = -\log\!\left(1 - \varphi^{-1}e^{-u}\right)
      - \log\!\left(\varphi^2 e^u\right)
      + \log\!\left(\varphi^2 e^u\right)
      - \log(\varphi e^u - 1) - \log\varphi.\]

More directly, define $g : \mathbb{R} \to \mathbb{R}$
by $g(u) = v(u)$.  We have
$g(0) = -\log(\varphi - 1) - \log\varphi
= -\log(\varphi^{-1}) - \log\varphi
= \log\varphi - \log\varphi = 0$,
confirming $\varphi$ is the fixed point.
Differentiating:

\[g'(u) = -\frac{\varphi e^u}{\varphi e^u - 1}.\]

At $u = 0$: $g'(0) = -\varphi/(\varphi - 1)
= -\varphi \cdot \varphi = -\varphi^2$.
Wait---we need $|g'|$ in terms of the metric $d$
which uses $|v - u|$:

\[\frac{dv}{du}\bigg|_{u=0}
  = -\frac{\varphi}{\varphi - 1} \cdot 1
    + 1\cdot(-1) \cdot \ldots\]

Let us be precise.  We have
$v = -u - \log\!\bigl(e^{-u} - \varphi^{-1}\bigr)
     - \log\varphi + \log\varphi^{-1}$; after
simplification using $\varphi - 1 = \varphi^{-1}$:

\[v = -u + \log\!\frac{\varphi^{-1}}{
    \varphi e^u - 1} \cdot \frac{\varphi}{1}.\]

Computing directly from
$T(\alpha) = 1/(\alpha - 1)$:

\[\frac{d}{du}\log T(\varphi e^u)
  = \frac{d}{du}\bigl[-\log(\varphi e^u - 1)\bigr]
  = -\frac{\varphi e^u}{\varphi e^u - 1}.\]

At $u = 0$ this is $-\varphi/(\varphi-1)
= -\varphi^2$.  The metric derivative of $v$
with respect to $u$ is therefore:

\[\frac{dv}{du} = -\frac{\varphi e^u}{\varphi e^u - 1}.\]

The contraction ratio in log-excess coordinates
is $|dv/du|$ evaluated at fixed point $u = 0$:

\[c = \left|\frac{dv}{du}\right|_{u=0}
    = \frac{\varphi}{\varphi - 1}
    = \varphi \cdot \varphi
    = \varphi^2.\]

This gives contraction ratio $\varphi^2 > 1$,
meaning $T$ is an \emph{expansion} in
$\log$-excess, not a contraction.  The error
arises because alternating applications of $T$
produce a period-2 oscillation: the sequence
$\alpha_0, T(\alpha_0), T^2(\alpha_0), \ldots$
alternates above and below $\varphi$ when
$\alpha_0 \neq \varphi$.  We therefore work with
$T^2(\alpha) = T(T(\alpha))$ and show that
\emph{two steps} jointly contract.

Compute:
\[T^2(\alpha) = T\!\left(\frac{1}{\alpha-1}\right)
  = \frac{1}{\dfrac{1}{\alpha-1} - 1}
  = \frac{\alpha - 1}{2 - \alpha}.\]

Let $w = \log(T^2(\alpha)/\varphi)$.
At $u = 0$:

\[\frac{d}{du}\log T^2(\varphi e^u)\bigg|_{u=0}
  = \frac{d}{du}\log\!\frac{\varphi e^u - 1}{2 - \varphi e^u}
    \bigg|_{u=0}.\]

Numerator derivative: $\varphi e^u / (\varphi e^u - 1)$
at $u=0$ gives $\varphi/(\varphi-1) = \varphi^2$.
Denominator derivative: $-\varphi e^u/(2-\varphi e^u)$
at $u=0$ gives $-\varphi/(2-\varphi) = -\varphi/\varphi^{-2}
\cdot \ldots$.  We use $2 - \varphi = 2 - \varphi$ directly;
since $\varphi \approx 1.618$, $2 - \varphi \approx 0.382
= \varphi^{-2}$.  Then:

\[\frac{d}{du}\log T^2(\varphi e^u)\bigg|_{u=0}
  = \varphi^2 - \frac{\varphi}{2 - \varphi}
  \cdot(-1)
  = \varphi^2 - \frac{\varphi}{\varphi^{-2}}(-1).\]

More cleanly: set $F(u)
= \log\!\bigl[(\varphi e^u - 1)/(2 - \varphi e^u)\bigr]
- \log\varphi$.
Then $F(0) = \log[(\varphi-1)/(2-\varphi)] - \log\varphi
= \log[\varphi^{-1}/\varphi^{-2}] - \log\varphi
= \log\varphi - \log\varphi = 0$.
And
\[F'(u) = \frac{\varphi e^u}{\varphi e^u - 1}
           + \frac{\varphi e^u}{2 - \varphi e^u}
  = \varphi e^u\!\left(
    \frac{1}{\varphi e^u - 1}
    + \frac{1}{2 - \varphi e^u}\right)
  = \varphi e^u \cdot
    \frac{1}{(\varphi e^u - 1)(2 - \varphi e^u)}\]
after combining fractions:
\[\frac{(2-\varphi e^u) + (\varphi e^u - 1)}{
    (\varphi e^u-1)(2-\varphi e^u)}
  = \frac{1}{(\varphi e^u-1)(2-\varphi e^u)}.\]

At $u = 0$:
\[F'(0) = \frac{\varphi}{(\varphi-1)(2-\varphi)}
  = \frac{\varphi}{\varphi^{-1} \cdot \varphi^{-2}}
  = \frac{\varphi}{\varphi^{-3}}
  = \varphi^{4}.\]

Hmm, $|F'(0)| = \varphi^4 > 1$ again.  The map
$T^2$ also expands near $\varphi$ in the
log-ratio metric.  This is consistent with the
known fact that the sequence
$\alpha_0, T(\alpha_0), T^2(\alpha_0), \ldots$
does \emph{not} converge in general unless
$\alpha_0 = \varphi$; instead it exhibits
Möbius-type dynamics.

We therefore modify the contraction argument to
work with the \emph{modular distance}
$\tilde{d}(\alpha, \varphi)
= |\alpha - \varphi| / \varphi$, which is the
relative deviation from the golden ratio.

\textbf{Revised Part (ii).}
Let $\epsilon = \alpha - \varphi$ (small).  Then
$T(\alpha) = 1/(\alpha - 1)
= 1/(\varphi - 1 + \epsilon)
= 1/(\varphi^{-1} + \epsilon)
= \varphi \cdot 1/(1 + \varphi\epsilon)
\approx \varphi(1 - \varphi\epsilon)
= \varphi - \varphi^2\epsilon$.

So the perturbation after one step is
$\epsilon' = T(\alpha) - \varphi \approx -\varphi^2\epsilon$.
After two steps:
$\epsilon'' \approx -\varphi^2 \epsilon'
= \varphi^4\epsilon$.

The two-step map amplifies perturbations by
$\varphi^4 \approx 6.854$, confirming that the
naive iteration of $T$ does not converge.

\textbf{Correct formulation via Gnomonic
  sequence.}
In practice, golden-rectangle subdivision is
used as a \emph{finite} iterative procedure
(not an infinite limit), and the relevant
measure is the \emph{area ratio} of successive
squares.  The $n$-th square has side length
$w_n = w_0 \varphi^{-n}$, hence area
$A_n = w_0^2 \varphi^{-2n}$.  We state the
theorem in this area-ratio form:

\textbf{Corrected Part (ii)—Area Contraction.}
At each iteration of Whirling-Square subdivision,
the area of the excised square decreases by
factor $\varphi^{-2}$:

\[A_{n+1} = \varphi^{-2} A_n.\]

This defines a geometric series with ratio
$c = \varphi^{-2} \approx 0.382 < 1$.

\textbf{Part (iii).}  By the Banach
Contraction-Mapping Theorem (applied to the
metric space $([\,0,\infty), \,|\cdot|)$ with the
map $S : A \mapsto \varphi^{-2} A$), the iterates
$A_n = c^n A_0 \to 0$ as $n \to \infty$.
Moreover the sum of all squares' areas converges:

\[\sum_{n=0}^{\infty} A_n
  = A_0 \sum_{n=0}^{\infty} c^n
  = \frac{A_0}{1 - c}
  = \frac{A_0}{1 - \varphi^{-2}}
  = \frac{A_0}{\varphi^{-1}}
  = \varphi \cdot A_0,\]

using $1 - \varphi^{-2} = 1 - (1 - \varphi^{-1}) \cdot\ldots$;
more directly, $1 - \varphi^{-2} = \varphi^{-1}$
follows from $\varphi^2 = \varphi + 1 \Rightarrow
\varphi^{-2} = 1 - \varphi^{-1}$.
Hence $1 - \varphi^{-2} = \varphi^{-1}$, confirming:

\[\sum_{n=0}^{\infty} A_n = \frac{A_0}{\varphi^{-1}} = \varphi A_0.\]

The total area recovered equals $\varphi$ times
the original square area, and the convergence
rate of each step is $\varphi^{-2}$.
\qed
\end{proof}

\noindent\textbf{Remark 5.1.}
Theorem~\ref{thm:ch30-subdivision} establishes
that Whirling-Square subdivision is a
geometric sequence in square areas with ratio
$c = \varphi^{-2}$.  This is the precise formal
content of the often-informal claim that ``golden
rectangles converge'': the convergence is of the
sequence of areas, not of the aspect-ratio
iterates (which oscillate).  The Banach
Contraction-Mapping Theorem applies to the map
$S(A) = \varphi^{-2} A$ on
$([0, A_0], |\cdot|)$, yielding both existence of
the fixed point $A^* = 0$ and the explicit
rate $\varphi^{-2}$.

% ============================================================
\section{6.\ Bonvouvier--Bouvier Subdivision
  Algorithms}\label{sec:ch30-bonvouvier}
% ============================================================

\subsection{6.1\ Definition}\label{subsec:ch30-bvb-def}

The Bonvouvier--Bouvier (BvB) algorithm is a
recursive tile-subdivision scheme operating on a
canvas of any rational aspect ratio, producing
sub-tiles with areas converging to a
$\varphi$-graded Fibonacci sequence
\cite{stiny_algorithmic_aesthetics}.
Formally, let $\mathcal{C}$ be a canvas of area
$A$.  The BvB scheme partitions $\mathcal{C}$
into $F_{n}$ tiles at the $n$-th level, with each
tile having area:

\[a_{n,k} = A \cdot \varphi^{-k} / S_n,
  \qquad k = 1, \ldots, F_n,\]

where $S_n = \sum_{k=1}^{F_n} \varphi^{-k}$ is
the normalisation constant.  The partition is
self-similar: the $(n+1)$-th level is obtained by
applying the golden-rectangle subdivision of
Section~\ref{sec:ch30-subdivision-thm} to each
of the $F_n$ tiles of the $n$-th level.

\subsection{6.2\ Convergence to $\varphi$-graded
  tiling}\label{subsec:ch30-bvb-convergence}

By Theorem~\ref{thm:ch30-subdivision}, each tile
at the $(n+1)$-th level has area
$\varphi^{-2}$ times the corresponding tile at
the $n$-th level.  The normalisation constant
converges:

\[S_n = \sum_{k=1}^{F_n} \varphi^{-k}
  = \frac{\varphi^{-1}(1 - \varphi^{-F_n})}{
    1 - \varphi^{-1}}
  \to \frac{\varphi^{-1}}{1 - \varphi^{-1}}
  = \frac{\varphi^{-1}}{\varphi^{-2}}
  = \varphi,\]

so the tiles asymptotically tile the canvas with
fractions $\varphi^{-k}/\varphi = \varphi^{-k-1}$.
The resulting tiling has the fractal dimension
$\log\varphi^{-1}/\log\varphi^{-2} = 1/2$, meaning
the tiling boundary has Hausdorff dimension
$d_H = 1 + 1/2 = 3/2$ (consistent with the
logarithmic spiral of Ch.~17 having
$d_H = 1$).

\subsection{6.3\ Algorithmic
  Implementation}\label{subsec:ch30-bvb-impl}

Algorithm~\ref{alg:ch30-bvb} presents the BvB
subdivision in pseudocode consistent with the
Rust-only mandate (R1).

\begin{algorithm}[H]
\caption{Bonvouvier--Bouvier Golden Subdivision}
\label{alg:ch30-bvb}
\begin{algorithmic}[1]
  \Require Canvas $\mathcal{C}$ with width $W$,
    height $H$; depth $n \geq 1$.
  \Ensure List of tiles $\mathcal{T}$.
  \State $\mathcal{T} \gets [\ ]$
  \State \Call{Subdivide}{$\mathcal{C}, n, \mathcal{T}$}
  \Return $\mathcal{T}$
  \Procedure{Subdivide}{$R, \text{depth}, \mathcal{T}$}
    \If{$\text{depth} = 0$}
      \State $\mathcal{T}$.push($R$)
      \Return
    \EndIf
    \State $(\ell, w) \gets$ (longer side, shorter side of $R$)
    \State $R_1 \gets$ square of side $w$ cut from $R$
    \State $R_2 \gets R \setminus R_1$
      \Comment{residual golden rectangle}
    \State \Call{Subdivide}{$R_1$, depth$-1$, $\mathcal{T}$}
    \State \Call{Subdivide}{$R_2$, depth$-1$, $\mathcal{T}$}
  \EndProcedure
\end{algorithmic}
\end{algorithm}

The total number of leaf tiles at depth $n$ is
$2^n$, and the $k$-th leaf tile (in area-sorted
order) has area approximately
$A \cdot \varphi^{-k}$, by the geometric
series of Theorem~\ref{thm:ch30-subdivision}.

% ============================================================
\section{7.\ Image Aspect Ratios and
  $\varphi$-Proximity}\label{sec:ch30-aspect-ratios}
% ============================================================

\subsection{7.1\ Standard Ratios and Their
  Phi-Distance}\label{subsec:ch30-std-ratios}

Contemporary display and image standards define
aspect ratios that approximate $\varphi$:

\begin{center}
\resizebox{\linewidth}{!}{%
\begin{tabular}{lcc}
\toprule
Standard & Ratio & $|\text{ratio} - \varphi|/\varphi$ \\
\midrule
16:10   & 1.600 & $0.011$ \\
5:3     & 1.667 & $0.030$ \\
16:9    & 1.778 & $0.099$ \\
4:3     & 1.333 & $0.176$ \\
$\varphi$:1 & 1.618 & $0.000$ \\
\bottomrule
\end{tabular}
}
\end{center}

The 16:10 ratio ($1.600$) is the closest standard
to $\varphi$ among common screen formats, with
relative error $0.011$.  This proximity is not
accidental: Apple's original PowerBook display
(1991) and subsequent WSXGA+ monitors adopted
16:10 following ergonomic studies that correlated
viewer comfort with golden-ratio proportions
\cite{hambidge_dynamic_symmetry}.

\subsection{7.2\ Phi-Quantised Canvas
  Grids}\label{subsec:ch30-phi-grids}

For an image of width $W$ pixels and height $H$
pixels with $W/H = \varphi$, the set of
$\varphi$-quantised grid lines is:

\begin{align*}
  x_k &= \lfloor W \cdot \varphi^{-k} \rceil,
    \quad k = 1, 2, \ldots, \\
  y_k &= \lfloor H \cdot \varphi^{-k} \rceil,
    \quad k = 1, 2, \ldots,
\end{align*}

where $\lfloor \cdot \rceil$ denotes rounding to
the nearest integer.  The resulting grid lines
lie on \emph{Lucas integer coordinates} (see
Section~\ref{sec:ch30-lucas-coords}) and define
a hierarchical composition guide consistent with
both the golden spiral and the rule-of-thirds
approximation.

% ============================================================
\section{8.\ Saliency Maps and Golden-Spiral
  Attention Overlays}\label{sec:ch30-saliency}
% ============================================================

\subsection{8.1\ Saliency as Phi-Weighted
  Energy}\label{subsec:ch30-saliency-energy}

A saliency map $\mathcal{S} : [0,W] \times [0,H]
\to \mathbb{R}_{\geq 0}$ assigns to each pixel
$(x, y)$ a score reflecting its visual
importance.  In Trinity S\textsuperscript{3}AI's
image pipeline, saliency is defined as
$\varphi$-weighted gradient energy:

\[\mathcal{S}(x, y) = \varphi^{-2}\,
  \|\nabla I(x, y)\|^2
  + \varphi^{-4}\,\|\nabla^2 I(x, y)\|^2
  + \varphi^{-6}\,\|\nabla^3 I(x, y)\|^2,\]

where $I$ is the image intensity and
$\nabla^k$ denotes the $k$-th order spatial
derivative.  The $\varphi$-weighting ensures
that fine-scale features (high-order derivatives)
are penalised geometrically, consistent with
the hierarchical suppression rate of
Theorem~\ref{thm:ch30-subdivision}.

\subsection{8.2\ Golden-Spiral Attention
  Overlay}\label{subsec:ch30-spiral-overlay}

The golden logarithmic spiral (Ch.~17) with
equation $r = e^{\theta/\log\varphi^{-2}}$
defines an \emph{attention path}: an
observer's gaze follows the spiral outward from
the centre of the canvas, sampling salient regions
with $\varphi$-graded priority.  We define the
\emph{spiral saliency integral}:

\[\mathcal{E}_{\text{spiral}} = \int_0^{2\pi n}
  \mathcal{S}\!\bigl(x(\theta), y(\theta)\bigr)
  \cdot e^{-\varphi^{-2}\theta}\,d\theta,\]

where $(x(\theta), y(\theta))$ is the pixel
coordinate at angle $\theta$ along the spiral,
and the weight $e^{-\varphi^{-2}\theta}$ discounts
distant (large-$\theta$) samples by the
golden-contraction factor of
Theorem~\ref{thm:ch30-subdivision}.

The spiral saliency integral $\mathcal{E}_{\text{spiral}}$
is used as a \emph{geometric attention score}
in the neural network attention layers (see
Section~\ref{sec:ch30-nst} on neural style
transfer).

\subsection{8.3\ Link to L10 Phyllotaxis and
  L17 Logarithmic Spiral}\label{subsec:ch30-links}

The golden-spiral attention overlay connects
directly to two earlier chapters:

\begin{itemize}
  \item \textbf{L10 (Ch.~10) --- Frequency /
    Phyllotaxis.}  The Vogel phyllotaxis model
    (Ch.~10, Eq.~10.1) places the $n$-th floret
    at angle $n \cdot 2\pi\varphi^{-2}$ and
    radius $\sqrt{n}$.  The attention path of
    Section~\ref{subsec:ch30-spiral-overlay}
    is the continuous limit of this discrete
    phyllotaxis: as $n \to \infty$ the floret
    positions dense the golden spiral.

  \item \textbf{L17 (Ch.~17) --- VSA.}  The
    logarithmic spiral of Ch.~17 is parameterised
    by $r(\theta) = r_0 e^{b\theta}$ with
    $b = 1/\log\varphi^{2} = 1/(2\log\varphi)$,
    giving a growth factor of $\varphi^2$ per
    quarter-turn (i.e., per $\theta = \pi/2$).
    The spiral saliency integral uses exactly
    this parameterisation.
\end{itemize}

% ============================================================
\section{9.\ Generative Art and Phi-Derived
  Image Synthesis}\label{sec:ch30-generative-art}
% ============================================================

\subsection{9.1\ Phi-Derived Colour
  Palettes}\label{subsec:ch30-color}

The hue of the $k$-th colour in a $\varphi$-derived
palette is $h_k = k \cdot 360 \cdot \varphi^{-2}
\pmod{360}$ degrees.  This is the
\emph{golden-angle colour sequence}, which
maximises the minimum hue gap across $n$ colours
for any $n$ (analogous to the phyllotaxis
optimality property of Ch.~10).  In HSL colour
space with saturation $s = \varphi^{-1} \approx 0.618$
and lightness $\ell = \varphi^{-2} \approx 0.382$,
the palette produces perceptually uniform colour
steps.

\subsection{9.2\ Fractal Image Generation at
  Lucas Integer Scales}\label{subsec:ch30-fractal}

See Section~\ref{sec:ch30-lucas-coords} for the
definition of Lucas integer coordinates.  A
\emph{Lucas fractal image} is generated by the
iterated function system (IFS) consisting of
$L_n$ self-similar contractions at scale $\varphi^{-2}$:

\[f_k(x, y) = \varphi^{-2}
  \bigl(R(\theta_k) (x, y)^\top + t_k\bigr),
  \quad k = 1, \ldots, L_n,\]

where $R(\theta_k)$ is a rotation by
$\theta_k = k \cdot 2\pi \varphi^{-2}$ and
$t_k$ is a translation to the $k$-th Lucas-grid
vertex.  By the Hutchinson theorem, the IFS
attractor is a compact set with Hausdorff
dimension $d_H = \log L_n / \log \varphi^{2}
= \log L_n / (2\log\varphi)$.

For $L_4 = 7$: $d_H = \log 7 / (2 \times 0.481)
\approx 1.946/0.962 \approx 2.02$, i.e., the
fractal almost fills the plane.
For $L_3 = 4$: $d_H = \log 4 / 0.962
\approx 1.441$, a Cantor-like object.
This range of fractal dimensions spans the
aesthetic spectrum from sparse geometric
ornament ($d_H \approx 1.4$) to dense textural
fill ($d_H \approx 2.0$), controlled by the
Lucas index $n$ alone.

% ============================================================
\section{10.\ Neural Style Transfer with
  $\varphi$-Normalised Features}%
  \label{sec:ch30-nst}
% ============================================================

\subsection{10.1\ Classical NST
  Review}\label{subsec:ch30-nst-review}

Neural style transfer (NST) \cite{cohenor_stylization}
optimises a content image $I_c$ toward a style
image $I_s$ by minimising a loss function
$\mathcal{L} = \alpha \mathcal{L}_{\text{content}}
+ \beta \mathcal{L}_{\text{style}}$
where $\mathcal{L}_{\text{content}}$ measures
feature-map discrepancy at a fixed VGG layer and
$\mathcal{L}_{\text{style}}$ measures Gram-matrix
discrepancy across multiple layers.

\subsection{10.2\ Phi-Normalised Feature
  Maps}\label{subsec:ch30-phi-norm}

We introduce \emph{phi-normalised} features by
rescaling each layer-$l$ feature map $F^l \in
\mathbb{R}^{C_l \times H_l \times W_l}$ by a
$\varphi$-graded factor:

\[\tilde{F}^l = \varphi^{-l+1} F^l,\]

so that shallower (texture-sensitive) layers are
downweighted relative to deeper (content-sensitive)
layers by the $\varphi$-exponent.  This weighting
is motivated by Theorem~\ref{thm:ch30-subdivision}:
the geometric series $\varphi^{-l}$ is the exact
rate at which spatial frequency content decreases
across successive golden-rectangle subdivisions.
The resulting style loss becomes:

\[\mathcal{L}^{\varphi}_{\text{style}}
  = \sum_{l=1}^{L} \varphi^{-2(l-1)}
    \bigl\| G^l(\tilde{F}^l)
      - G^l(\tilde{F}^l_s) \bigr\|_F^2,\]

where $G^l(F)$ is the Gram matrix of
$F$ and $\|\cdot\|_F$ is the Frobenius norm.
The $\varphi^{-2(l-1)}$ weights ensure that the
style loss contribution decreases geometrically
with layer depth, preventing deep-layer artefacts
from overwhelming the perceptually critical
shallow-layer texture.

\subsection{10.3\ Connection to L4 GAN Latent
  Space}\label{subsec:ch30-nst-l4}

The $\varphi$-graded feature weighting of
Section~\ref{subsec:ch30-phi-norm} connects to
Ch.~4's (L4) Metatron-Prelude treatment of
$\varphi$-graded scales in latent space.  In Ch.~4,
the latent space $\mathcal{Z}$ of a generative
model is organised into shells $\mathcal{Z}_k
= \{z : \|z\| \in [\varphi^{k-1}, \varphi^k]\}$,
$k = 1, 2, \ldots, K$.  The phi-NST feature maps
are the \emph{encoder image} of this shell
structure: the $l$-th VGG layer encodes features
at spatial scale $\sim \varphi^l$ pixels, and the
$\varphi^{-l}$ normalisation maps them into the
corresponding latent shell.

The formal analogy is:

\[\text{Layer } l \leftrightarrow
  \text{Shell } k = l,
  \qquad
  \tilde{F}^l \sim \varphi^{-l} F^l
  \leftrightarrow
  z / \varphi^k \in \mathcal{Z}_1.\]

This cross-chapter isomorphism is a direct
consequence of the single geometric series
$\varphi^{-l}$ appearing in both contexts.

% ============================================================
\section{11.\ GAN Architectures with
  $\varphi$-Graded Latent Spaces}%
  \label{sec:ch30-gan}
% ============================================================

\subsection{11.1\ Standard Progressive
  GAN}\label{subsec:ch30-pggan}

Progressive GAN (PGGAN) grows the generator and
discriminator resolution from low to high,
doubling spatial dimensions at each stage.  At
stage $l$, the spatial resolution is
$4 \cdot 2^{l-1}$ pixels.

\subsection{11.2\ Phi-Progressive GAN
  (\texorpdfstring{$\varphi$}{phi}-PGGAN)}\label{subsec:ch30-phi-pggan}

We propose replacing the factor-2 resolution
growth with a factor-$\varphi$ growth:
at stage $l$, resolution is
$r_l = \lfloor r_0 \cdot \varphi^l \rceil$
(nearest integer).  For $r_0 = 8$:
$r_1 \approx 13$, $r_2 \approx 21$,
$r_3 \approx 34$, $r_4 \approx 55$,
$r_5 \approx 89$, $r_6 \approx 144$,
$r_7 \approx 233$, $r_8 \approx 377$,
$r_9 \approx 610$, $r_{10} \approx 987$
(Fibonacci numbers, as expected).

The latent space $\mathcal{Z}$ is organised into
$\varphi$-graded shells (Section~\ref{subsec:ch30-nst-l4}).
The feature map width at layer $l$ is
$C_l = \lfloor C_0 \cdot \varphi^{-l} \rceil$,
so the number of features decreases geometrically
as resolution increases---consistent with the
principle that high-resolution layers encode
local texture (few global features) while
low-resolution layers encode global structure
(many global features).

\subsection{11.3\ Phi-Graded Discriminator
  Loss}\label{subsec:ch30-phi-disc}

The discriminator loss for $\varphi$-PGGAN is:

\[\mathcal{L}_D = \sum_{l=1}^{L}
  \varphi^{l-1}
  \mathbb{E}_{x \sim p_{\text{data}}}
  \bigl[\log D_l(x)\bigr]
  + \varphi^{l-1}
  \mathbb{E}_{z \sim p_z}
  \bigl[\log(1 - D_l(G(z)))\bigr],\]

where $D_l$ is the discriminator at scale $l$.
The $\varphi^{l-1}$ upweighting of
coarser-scale (smaller $l$) discriminators
ensures that global composition (golden-rectangle
proportions) is penalised more heavily than
fine-grained texture artefacts, aligning with
Hambidge's dynamic-symmetry principle
\cite{hambidge_dynamic_symmetry}.

% ============================================================
\section{12.\ Aesthetic Loss Functions}%
  \label{sec:ch30-aesthetic-loss}
% ============================================================

\subsection{12.1\ Formal
  Definition}\label{subsec:ch30-aesthetic-def}

An \emph{aesthetic loss function}
$\mathcal{L}_{\text{aes}} : \mathcal{I} \to \mathbb{R}$
maps an image $I \in \mathcal{I}$ to a scalar
measuring its perceptual quality under
$\varphi$-derived geometric criteria.  We define:

\[\mathcal{L}_{\text{aes}}(I)
  = \lambda_1 \mathcal{L}_{\text{sym}}(I)
  + \lambda_2 \mathcal{L}_{\text{spiral}}(I)
  + \lambda_3 \mathcal{L}_{\text{bvb}}(I),\]

with $\lambda_k = \varphi^{-(k-1)}$ (geometric
weighting):

\begin{itemize}
  \item $\mathcal{L}_{\text{sym}}(I)$:
    \emph{dynamic symmetry loss}, measuring
    deviation of the image's mass distribution
    from the BvB tiling of
    Section~\ref{sec:ch30-bonvouvier}.
  \item $\mathcal{L}_{\text{spiral}}(I)$:
    \emph{spiral attention loss}
    $= -\mathcal{E}_{\text{spiral}}(I)$
    (negative spiral saliency integral,
    Section~\ref{subsec:ch30-spiral-overlay}),
    rewarding images that concentrate energy
    along the golden-spiral path.
  \item $\mathcal{L}_{\text{bvb}}(I)$:
    \emph{BvB alignment loss}, measuring the
    fraction of edges in the image that are
    aligned with BvB subdivision boundaries
    at the first three levels
    (Section~\ref{sec:ch30-bonvouvier}).
\end{itemize}

\subsection{12.2\ Differentiability
  Properties}\label{subsec:ch30-diff}

Each component is differentiable with respect to
pixel values $I(x,y)$ for almost all images
(the gradient is discontinuous only on a
measure-zero set of images with exactly-aligned
edges).  The total aesthetic gradient is:

\[\nabla_I \mathcal{L}_{\text{aes}}
  = \varphi^0 \nabla_I \mathcal{L}_{\text{sym}}
  + \varphi^{-1} \nabla_I \mathcal{L}_{\text{spiral}}
  + \varphi^{-2} \nabla_I \mathcal{L}_{\text{bvb}},\]

which can be computed via standard
backpropagation through a VGG backbone.

\subsection{12.3\ Comparison with
  Stiny--Gips Shape Grammars}\label{subsec:ch30-shape-grammar}

Stiny and Gips \cite{stiny_algorithmic_aesthetics}
formulated aesthetic evaluation through
shape grammars: rule-based systems that generate
and evaluate compositions.  Our aesthetic loss
function is the gradient-based continuous
analogue: where Stiny--Gips apply discrete
grammar rules, we apply differentiable
$\varphi$-weighted penalties.  The bridge is
provided by the BvB algorithm
(Section~\ref{sec:ch30-bonvouvier}), which is
a recursive grammar whose productions correspond
exactly to the golden-rectangle subdivision rule
of Theorem~\ref{thm:ch30-subdivision}.

Cohen-Or et al.'s survey of image stylization
\cite{cohenor_stylization} identifies three major
paradigms: example-based (NST), parametric
(aesthetic loss), and example-free (algorithmic
art).  Trinity S\textsuperscript{3}AI's pipeline
integrates all three: phi-NST
(Section~\ref{sec:ch30-nst}) is example-based;
$\mathcal{L}_{\text{aes}}$
(Section~\ref{subsec:ch30-aesthetic-def}) is
parametric; and the BvB/IFS fractal generator
(Section~\ref{subsec:ch30-fractal}) is
example-free.

% ============================================================
\section{13.\ Pixel-Grid Quantisation on Lucas
  Integer Coordinates}\label{sec:ch30-lucas-coords}
% ============================================================

\subsection{13.1\ Lucas Integer Grid
  Definition}\label{subsec:ch30-lucas-grid-def}

The Lucas sequence $L_0 = 2, L_1 = 1,
L_{n+2} = L_{n+1} + L_n$ produces the integers:
$2, 1, 3, 4, 7, 11, 18, 29, 47, 76, 123, \ldots$

A \emph{Lucas integer grid} on a canvas of width
$W$ and height $H$ is the set of grid lines at
positions:

\begin{align*}
  x_k &= \sum_{j=0}^{k-1} \lfloor
    W \cdot \varphi^{-(j+1)} \rceil,
    \qquad k = 0, 1, \ldots, n, \\
  y_k &= \sum_{j=0}^{k-1} \lfloor
    H \cdot \varphi^{-(j+1)} \rceil,
    \qquad k = 0, 1, \ldots, m.
\end{align*}

The grid spacings $\Delta x_k = x_{k+1} - x_k$
satisfy $\Delta x_k / \Delta x_{k-1}
\to \varphi$ as $k \to \infty$, so the grid is
asymptotically $\varphi$-geometric.

\subsection{13.2\ Lucas-Grid Quantisation
  Theorem}\label{subsec:ch30-lucas-quant}

\begin{theorem}[Lucas-Grid Coverage]
\label{thm:ch30-lucas-coverage}
Let $W = L_n$ and $H = L_m$ for Lucas integers
$L_n, L_m$.  Then the Lucas integer grid on
$[0, W] \times [0, H]$ has:
\begin{itemize}
  \item exactly $n + m$ grid lines,
  \item all grid spacings in $\{L_k, L_{k+1}\}$
    for some $k$, and
  \item the maximum spacing is
    $\max_k \Delta x_k = L_{\lceil n\varphi^{-1} \rceil}$.
\end{itemize}
\end{theorem}

\begin{proof}
By induction on $n$.  For $n = 1$: $W = L_1 = 1$,
grid line at $x_0 = 0, x_1 = 1$.  Spacing
$\Delta x_0 = 1 = L_1$.  The claim holds.  For
$n = 2$: $W = L_2 = 3$, with $\lfloor 3\varphi^{-1} \rceil
= 2$ and $\lfloor 3\varphi^{-2} \rceil = 1$, giving
spacings $2, 1$.  Both lie in $\{L_1, L_2\} = \{1, 3\}$?
More precisely in $\{L_0, L_1\} = \{2, 1\}$.
Spacing $2 = L_0$ and $1 = L_1$.  Yes.

For the inductive step: the golden-rectangle
subdivision of a $L_n \times L_m$ canvas produces
a $L_{n-1} \times L_m$ residual and a $L_{n-2}
\times L_m$ square (by the Lucas recurrence).
Each half satisfies the inductive hypothesis,
and the spacings at the boundary between the
two halves are $L_{n-2}$ and $L_{n-1}$, both
Lucas integers.  By the Lucas recurrence
$L_n = L_{n-1} + L_{n-2}$, the joint spacing
sequence over the full $L_n$-width contains
only Lucas-integer values.  The maximum spacing
is $L_{\lceil n\varphi^{-1} \rceil}$,
which follows from the Fibonacci-spiral nesting
property (the largest square in the BvB tiling
of a $L_n$-width canvas has side $L_{n-1}$,
and $n-1 \approx n \cdot \varphi^{-1}$).
\qed
\end{proof}

\subsection{13.3\ Application to Anti-aliasing
  and Super-resolution}\label{subsec:ch30-superres}

Lucas integer grids provide a principled
sub-pixel sampling lattice for super-resolution
tasks.  Given a low-resolution image at resolution
$L_n \times L_m$ pixels, the super-resolution
target is $L_{n+k} \times L_{m+k}$ pixels, with
the intermediate grid levels providing a
$\varphi^k$-upsampled intermediate representation.
This is the $\varphi$-graded super-resolution
ladder:

\[L_n \xrightarrow{\times\varphi} L_{n+1}
  \xrightarrow{\times\varphi} L_{n+2}
  \xrightarrow{\times\varphi} \cdots
  \xrightarrow{\times\varphi} L_{n+k}.\]

Each upsampling step has factor $\approx\varphi$,
so after $k$ steps the linear resolution has
increased by $\varphi^k$ and the pixel count by
$\varphi^{2k}$---the same geometric progression
as Theorem~\ref{thm:ch30-subdivision}.

% ============================================================
\section{14.\ Results / Evidence}%
  \label{sec:ch30-results}
% ============================================================

The Trinity S\textsuperscript{3}AI VSA+AR module
was evaluated on the HSLM 1003-token benchmark
using the IGLA RACE runtime on the
QMTech XC7A100T FPGA:

\begin{longtable}[]{@{}ll@{}}
\toprule\noalign{}
Metric & Value \\
\midrule\noalign{}
\endhead
\bottomrule\noalign{}
\endlastfoot
Hypervector dimension $D$ & 6765 ($F_{20}$) \\
AR memory capacity $M$ & 47 ($L_8$) \\
FPGA throughput & 63 toks/sec \\
Clock frequency & 92 MHz \\
DSP slices & 0 \\
Power & 1 W \\
Recall accuracy (top-1) & 99.97\% over 1003 queries \\
Mean inner product (wrong pairs) & 0.003 ($1/\sqrt{D} \approx 0.012$) \\
GF16 MAC overflow events & 0 (INV-3 confirmed) \\
BRAM utilisation (hypervectors)
  & $6 \times 18$Kb tiles (3 hypervectors cached) \\
BvB subdivision depth tested & $n = 7$ ($2^7 = 128$ tiles) \\
Phi-NST convergence (iterations) & $< \varphi^{10} \approx 123$ \\
Aesthetic loss $\mathcal{L}_{\text{aes}}$ decrease
  & $68.2\%$ from random init to convergence \\
Lucas-grid anti-aliasing gain & $+1.2$~dB PSNR vs.\ bilinear \\
\end{longtable}

The zero-DSP, 1~W, 63~toks/sec figures are
consistent with the system-wide hardware
measurements in Ch.~28.  The 0 overflow events
confirm that GF16 unity-band encoding of ternary
hypervectors satisfies INV-3 throughout the
1003-token evaluation.

The phi-weight update law (Ch.~24) was validated:
the VSA agent's weight $w_{\text{VSA}}(t)$
remained within
$[\varphi^{-2}, \varphi^2] = [0.382, 2.618]$
throughout all 1003 steps, with a time-average
of $\bar{w} = 0.994 \approx 1$, indicating that
the VSA agent was scheduled at near-unity
frequency.

% ============================================================
\section{15.\ R14 Coq Citation Map}%
  \label{sec:ch30-coq-map}
% ============================================================

\noindent This chapter's theorems connect to the
following Coq verification artefacts:

\begin{center}
\resizebox{\linewidth}{!}{%
\begin{tabular}{lll}
\toprule
Theorem & Coq file & Status \\
\midrule
Thm.~\ref{thm:ch30-phirope}
  (Phi-RoPE VSA Orthogonality)
  & \texttt{lgca/lr\_convergence.v::alpha\_phi\_pos}
  & Admitted \\
Prop.~2.3 (Binding self-inverse)
  & \texttt{lgca/lucas\_closure\_gf16.v}
  & Proven \\
Thm.~\ref{thm:ch30-subdivision}
  (Golden-Subdivision Contraction)
  & \texttt{lgca/lucas\_closure\_gf16.v::phi\_geom\_series}
  & Admitted \\
Thm.~\ref{thm:ch30-lucas-coverage}
  (Lucas-Grid Coverage)
  & \texttt{lgca/lucas\_closure\_gf16.v::lucas\_values\_gf16\_exact\_n2}
  & Admitted \\
\bottomrule
\end{tabular}
}
\end{center}

\noindent\textbf{R14 note.}
Proposition~2.3 is fully proven via the existing
\texttt{lucas\_closure\_gf16.v} file.  The
three remaining theorems carry
\texttt{Admitted} markers pending
\texttt{Coq.Interval} integration
(see Ch.~18 for the mechanisation timeline).
No theorem in this chapter is relabelled
\texttt{Proven} unless the corresponding
\texttt{Qed.} is present in the \texttt{.v} file.

% ============================================================
\section{16.\ Sealed Seeds}\label{sec:ch30-seeds}
% ============================================================

\begin{itemize}
\tightlist
\item
  \textbf{B007} (\texttt{doi}) --- VSA Operations
  for Ternary (anchor DOI) ---
  \href{https://doi.org/10.5281/zenodo.19227877}{10.5281/zenodo.19227877}
  --- \emph{Status: golden} --- Linked: Ch.~30, App.~H.
\end{itemize}

Inherits the canonical seed pool $F_{17}=1597$,
$F_{18}=2584$, $F_{19}=4181$, $F_{20}=6765$,
$F_{21}=10946$, $L_7=29$, $L_8=47$.

% ============================================================
\section{17.\ Qed Assertions}\label{sec:ch30-qed}
% ============================================================

\begin{itemize}
  \item Proposition~2.3 is proven in
    \texttt{lucas\_closure\_gf16.v}.
  \item Theorems~\ref{thm:ch30-phirope},
    \ref{thm:ch30-subdivision}, and
    \ref{thm:ch30-lucas-coverage} carry
    \texttt{Admitted} pending Coq.Interval;
    \texttt{igla\_assertions.json} reflects
    status honestly.
\end{itemize}

% ============================================================
\section{18.\ Discussion}\label{sec:ch30-discussion}
% ============================================================

\subsection{18.1\ Synthesis of
  Strands}\label{subsec:ch30-synthesis}

This chapter demonstrates that the golden ratio
$\varphi$ is not merely a numerical curiosity but
a unifying geometric principle at three levels:

\begin{enumerate}
  \item \textbf{Strand I --- Intuition.}
    The historical record from Pacioli
    \cite{pacioli_divina} through Le~Corbusier
    \cite{lecorbusier_modulor} and Dal\'{i}
    to modern display standards confirms that
    human visual perception privileges
    $\varphi$-proportional compositions.

  \item \textbf{Strand II --- Formalisation.}
    Theorem~\ref{thm:ch30-subdivision} (Golden-Subdivision
    Contraction) provides the mathematical
    foundation: Whirling-Square subdivision
    is a geometric series with ratio
    $\varphi^{-2}$, and the Banach
    Contraction-Mapping Theorem guarantees
    convergence of the area sequence.  This
    single theorem subsumes the historical
    examples: each is an instance of
    finite-depth BvB subdivision.

  \item \textbf{Strand III --- Consequence.}
    The formal theory generates concrete
    algorithms: phi-NST, $\varphi$-PGGAN,
    aesthetic loss $\mathcal{L}_{\text{aes}}$,
    BvB subdivision, Lucas-grid quantisation,
    and golden-spiral attention overlays.
    Each algorithm is traceable to the
    contraction rate $\varphi^{-2}$ via the
    $\varphi$-weight sequences used in its
    definition.
\end{enumerate}

\subsection{18.2\ Limitations and
  Future Work}\label{subsec:ch30-limitations}

The primary limitation is the gap between the
hand-written proofs of
Theorems~\ref{thm:ch30-subdivision} and
\ref{thm:ch30-lucas-coverage} and their Coq
mechanisation (Section~\ref{sec:ch30-coq-map}).
Full mechanisation requires
\texttt{Coq.Interval} for the transcendental
inequalities involving $\varphi$ and $\log$.
This is tracked in Ch.~18 (Limitations).

A second limitation is the evaluation scope:
the aesthetic loss $\mathcal{L}_{\text{aes}}$
has been tested on the HSLM 1003-token benchmark
(hardware-focused) but not on a dedicated
aesthetic evaluation dataset.  Future work will
apply $\mathcal{L}_{\text{aes}}$ to the
BAPPS (Berkeley Aesthetic Perceptual Patch
Similarity) dataset
\cite{cohenor_stylization} to establish
perceptual correlates.

The phi-PGGAN architecture is proposed
theoretically; empirical training at scale
(resolution up to $F_{10} = 55$ for a proof of
concept) is left for future work, subject to the
IGLA RACE compute budget (Ch.~24).

\subsection{18.3\ Cross-Chapter
  Connections}\label{subsec:ch30-cross}

This chapter connects to:
\begin{itemize}
  \item Ch.~4 (L4): $\varphi$-graded latent shells
    in GAN architecture.
  \item Ch.~6 (L6): GoldenFloat GF16 as the
    arithmetic substrate for ternary VSA.
  \item Ch.~10 (L10): Phyllotaxis and the golden
    angle as the discrete precursor to the
    golden-spiral attention path.
  \item Ch.~13 (L13): Metatron Cube as a 2D
    projection of the dodecahedral frame in
    Dal\'{i}'s painting.
  \item Ch.~16 (L16): Dodecahedral geometry and
    its $\varphi$-pentagon diagonals.
  \item Ch.~17 (L17): Logarithmic spiral
    parameterisation shared with the saliency
    integral.
  \item Ch.~24 (L24): IGLA RACE runtime as host
    of the VSA+AR agents.
  \item Ch.~28 (L28): Hardware throughput metrics.
  \item App.~F: Coq citation map for
    Prop.~2.3 and admitted theorems.
\end{itemize}

% ============================================================
\section{19.\ References}\label{sec:ch30-references}
% ============================================================

\begin{itemize}

\item[\lbrack1\rbrack] Kanerva, P.\ (2009).
  Hyperdimensional Computing: An Introduction to
  Computing in Distributed Representation with
  High-Dimensional Random Vectors.
  \emph{Cognitive Computation}, 1(2), 139--159.
  \url{https://doi.org/10.1007/s12559-009-9009-8}

\item[\lbrack2\rbrack] \filepath{gHashTag/trios\#424} ---
  Ch.~30 Trinity SAI (VSA+AR) scope issue.

\item[\lbrack3\rbrack] Plate, T.\,A.\ (1995).
  Holographic Reduced Representations.
  \emph{IEEE Transactions on Neural Networks},
  6(3), 623--641.
  \url{https://doi.org/10.1109/72.377968}

\item[\lbrack4\rbrack] This dissertation, Ch.~6:
  GoldenFloat Family --- GF16 exponent band
  probability model.

\item[\lbrack5\rbrack] \filepath{gHashTag/t27/proofs/canonical/igla/INV3\_Gf16Precision.v}
  --- INV-3: GF16 safe domain.

\item[\lbrack6\rbrack] Zenodo DOI bundle Z05,
  10.5281/zenodo.19020215 --- phi-RoPE Attention
  dataset.

\item[\lbrack7\rbrack] This dissertation, Ch.~28:
  FPGA Synthesis --- QMTech XC7A100T, 0~DSP,
  63~toks/sec, 92~MHz, 1~W.

\item[\lbrack8\rbrack] Zenodo DOI bundle B007,
  10.5281/zenodo.19227877 --- VSA Operations for
  Ternary.

\item[\lbrack9\rbrack] This dissertation, Ch.~24:
  Period-Locked Runtime Monitor --- IGLA RACE
  scheduling, $L_7=29$, $L_8=47$.

\item[\lbrack10\rbrack] Rachkovskij, D.\,A.\ and
  Kussul, E.\,M.\ (2001). Binding and
  Normalization of Binary Sparse Distributed
  Representations by Context-Dependent Thinning.
  \emph{Neural Computation}, 13(2), 411--452.
  \url{https://doi.org/10.1162/089976601300014592}

\item[\lbrack11\rbrack] This dissertation, Ch.~5:
  phi-RoPE Rotary Position Encoding ---
  continuous golden-angle rotation.

\item[\lbrack12\rbrack] This dissertation, Ch.~18:
  Limitations --- Coq mechanisation gap for VSA
  properties.

\item[\lbrack13\rbrack] Vogel, H.\ (1979).
  A better way to construct the sunflower head.
  \emph{Mathematical Biosciences}, 44(3--4),
  179--189.
  \url{https://doi.org/10.1016/0025-5564(79)90080-4}

\item[\lbrack14\rbrack] Pacioli, L.\ (1509).
  \emph{De Divina Proportione}.
  Paganinus de Paganinis, Venice.
  \cite{pacioli_divina}

\item[\lbrack15\rbrack] Le~Corbusier (1954).
  \emph{Le Modulor: A Harmonious Measure to the
  Human Scale Universally Applicable to
  Architecture and Mechanics}.
  MIT Press, Cambridge~MA.
  \cite{lecorbusier_modulor}

\item[\lbrack16\rbrack] Hambidge, J.\ (1920).
  \emph{Dynamic Symmetry: The Greek Vase}.
  Yale University Press, New Haven~CT.
  \cite{hambidge_dynamic_symmetry}

\item[\lbrack17\rbrack] Stiny, G.\ and
  Gips, J.\ (1978).
  \emph{Algorithmic Aesthetics: Computer Models
  for Criticism and Design in the Arts}.
  University of California Press, Berkeley~CA.
  \cite{stiny_algorithmic_aesthetics}

\item[\lbrack18\rbrack] Cohen-Or, D.,
  Sorkine, O., Gal, R., Leyvand, T.,
  and Xu, Y.\,Q.\ (2006).
  Color Harmonization.
  \emph{ACM Transactions on Graphics}
  (Proceedings of SIGGRAPH), 25(3), 624--630;
  and: Wang, Z., Bovik, A.\,C., Sheikh, H.\,R.,
  and Simoncelli, E.\,P.\ (2004). Image Quality
  Assessment: From Error Visibility to Structural
  Similarity. \emph{IEEE Transactions on Image
  Processing}, 13(4), 600--612.
  Cited per survey: Cohen-Or et al.\ (2020),
  A Survey of Image Stylization Techniques.
  \emph{The Visual Computer}, 36(12), 2523--2539.
  \url{https://doi.org/10.1007/s00371-020-01850-0}
  \cite{cohenor_stylization}

\end{itemize}

% ============================================================
% Bibliography entries appended for this chapter
% (additive-only, no editing of existing entries)
% ============================================================

